\chapter{Results}
\label{ch:results}

\section{How to Read This Chapter}
\label{sec:results-reading}

Every number in this chapter is generated from a committed evaluation report rather than
transcribed, and every table states the corpus, model stack, numeric precision, scoring ruler
and aggregation it was produced under. Two tables that differ in any of those are not
comparable, and the protocol notes are what make that checkable instead of assumed.

Three conventions govern what may be claimed.

\textbf{In-sample figures are never reported as accuracy.} Some measurements in this project
were taken on documents that also informed prompt design, or on small development subsets, or
under scoring rules later found defective. Those results are diagnostic: they answer
\emph{which approach is worth pursuing}, not \emph{how well does the system perform}. They are
confined to \S\ref{sec:results-preliminary} and labelled there.

\textbf{Perfect-text figures are never reported as deliverable capability.} Where the
structuring stage is measured on text rendered from the reference data, the number is an upper
bound obtained with an instrument that does not exist at inference time
(\S\ref{sec:method-attribution}). It appears throughout as an attribution device and nowhere as
a claim about the system.

\textbf{Superseded figures are not quoted for precision.} Several measurements in this project
were re-scored after the corrections of Chapter~\ref{ch:measurement-validity}. Where an earlier
figure matters to an argument, this chapter refers to what changed and why, not treating
a superseded number as if it still stood.

\section{Preliminary Studies}
\label{sec:results-preliminary}

Three studies shaped the design before any sealed measurement was taken. All three are
in-sample, all three predate the scoring corrections, and none of their figures is an accuracy
claim. They are reported because the decisions they produced shaped the architecture, and a reader
entitled to question those decisions is entitled to see what informed them.

\textbf{A survey of local candidate readers} evaluated a cohort of open-weight document and
vision-language models on a small set of invoices under the local hardware envelope. Its
function was elimination, not ranking. It removed candidates that could not be run at all on
the target hardware, candidates whose available ports were incompatible with the platform, and
candidates that produced degenerate output. The scores it produced are low and are superseded
twice over --- by a change of reader and by the scoring corrections --- so they are not quoted.
What survived the survey is what mattered: a shortlist of readers worth measuring properly.

\textbf{A probe of structured-output prompting} asked whether instructing a reading model to
emit structured data directly, rather than transcribing to text for a second model to structure,
was competitive. The literature gave reason for caution: format restrictions have been shown to
degrade model performance on reasoning-adjacent tasks rather than being a free convenience
\parencite{tam2024format}. The probe's first verdict had to be withdrawn --- the measurement harness
was silently discarding valid model output, one model emitting a well-formed object per page
that the adapter could not parse once pages were concatenated. After repair, the finding stood
qualitatively: direct structured emission was not competitive with the two-stage arrangement on
this corpus, and the two-stage design was retained.

\textbf{A comparison of pipeline arrangements} measured a rule-based extraction baseline against
a single-shot arrangement and against the orchestrated two-stage arrangement on a small
development subset. The orchestrated arrangement won clearly, and that ordering is the reason
the architecture of Chapter~\ref{ch:system-design} has the shape it does. Two details from that
study are worth carrying forward. First, all three arms emitted \emph{no} values for absent
fields, which is early evidence that the honest-null contract is satisfiable, not merely
aspirational. Second, a token-budget defect was depressing the orchestrated arm until it was
found --- the model was being cut off before it finished emitting, and the resulting score was a
measurement of the budget, not of the approach.

\section{Reader Selection}
\label{sec:results-reader}

The pipeline's reading stage was selected by competitive comparison. This is worth stating
plainly because the structuring stage was not, and \S\ref{sec:results-adaptation} returns to
that asymmetry as a limitation.

\subsection{Why aggregate scores could not decide it}

Two candidates emerged from the survey as finalists: Qwen3-VL-4B-Instruct, a general-purpose
vision-language model of 4.4 billion parameters \parencite{qwen3vl}, and olmOCR-2-7B, a
purpose-built document-transcription model of 8.3 billion \parencite{olmocr2}.\footnote{Both
names cite the language backbone rather than the full multimodal checkpoint; the figures given
here are the checkpoints as published.} The two available
instruments told conflicting stories about them. On end-to-end score under the pre-correction
ruler, the transcription model led by two to five points (Table~\ref{tab:reader-lineage}). On
reading quality measured directly on the transcript (\S\ref{sec:method-answerability}), the
general-purpose model led narrowly under that same defective ruler (0.777 against 0.774) ---
and when the corrections of Chapter~\ref{ch:measurement-validity} were applied to the
identical stored transcripts, that ordering \emph{inverted} (0.906 against 0.913), both
readings within a point of each other.

An ordering that flips when the ruler is corrected is not evidence that the second model is
better. It is evidence that \emph{the instrument does not separate these two candidates}, and
that a selection made on it would have been arbitrary. That is the finding, and it forced the
decision onto the individual misses.

One asymmetry in the record must be stated. After the ruler repairs, only the selected reader
was re-measured end-to-end (the fourth row of Table~\ref{tab:reader-lineage}); the
transcription model's post-correction end-to-end score was never taken. Its transcripts are
retained, so the measurement remains possible --- but by that point the selection no longer
rested on end-to-end score at all, because an instrument that had already flipped its verdict
once under a ruler repair had disqualified itself as the deciding instrument.

% GENERATED FILE -- DO NOT EDIT BY HAND.
% Regenerate with: make thesis-assets  (scripts/thesis_assets.py)
% Editing this file by hand breaks the guarantee that every number in the
% manuscript is derived from committed evidence (ADR-055).

% Source(s): data/finetune/eval-zeroshot-val.json; data/finetune/eval-zeroshot-qwen-val.json; data/finetune/eval-zeroshot-olmocr-val.json; data/finetune/eval-zeroshot-qwen-adr059-val.json

\begin{table}[htbp]
  \centering
  \caption{End-to-end score by reading model, with the structurer and the corpus held fixed.}
  \label{tab:reader-lineage}
  \begin{tabular}{lr}
    \toprule
    Reader & Overall F$_1$ \\
    \midrule
    granite-docling-258M & 0.6771 \\
    Qwen3-VL-4B-Instruct & 0.7829 \\
    olmOCR-2-7B & 0.8335 \\
    \midrule
    Qwen3-VL-4B-Instruct \emph{(after ruler repairs)} & 0.8257 \\
    \bottomrule
  \end{tabular}

  \vspace{4pt}
  \begin{minipage}{0.95\linewidth}\footnotesize
  \textbf{Protocol note.} One structurer, one corpus, one prompt: the reader is the only thing that changes across the first three rows. The fourth row is the same reader as the second, re-scored after the scoring defects of Chapter~\ref{ch:measurement-validity} were repaired --- so the gap between rows two and four is instrument, not model. This is why reader selection could not be settled on end-to-end score alone: the ranking moved when the ruler was fixed, which is evidence about the ruler and not about the readers.
  \end{minipage}
\end{table}


\subsection{The per-miss audit, and its verdict}

The decision therefore moved to the individual misses. Every remaining miss of the
then-leading transcription model --- 52 individual ``missing'' values across the 29 sealed
validation invoices --- was judged by hand against the page image and against both finalists'
transcripts, and classified as a genuine reading failure, a comparison-rule artefact, or a
value not findable on the page by anything. That audit produced 23 exclusions --- values
absent from the page and from its own embedded text layer, which no vision system could
recover --- plus six ruler repairs. A symmetric pass over the general-purpose model's own
residual misses followed, exposing two further ruler gaps, so that neither finalist was
audited under a rule the other was spared.
Table~\ref{tab:reader-findability} reports the corrected result.

% GENERATED FILE -- DO NOT EDIT BY HAND.
% Regenerate with: make thesis-assets  (scripts/thesis_assets.py)
% Editing this file by hand breaks the guarantee that every number in the
% manuscript is derived from committed evidence (ADR-055).

% Source(s): scripts/findability_corrected.py (over data/finetune/bakeoff/)

\begin{table}[htbp]
  \centering
  \caption{Reading quality per candidate after every miss was judged by hand against the page image.}
  \label{tab:reader-findability}
  \begin{tabular}{lrr}
    \toprule
    Reader & Corrected findability & Real misses \\
    \midrule
    Embedded text layer \emph{(ceiling)} & 0.995 & 3 \\
    Qwen3-VL-8B-Instruct & 0.976 & 13 \\
    Qwen3-VL-4B-Instruct & 0.970 & 16 \\
    olmOCR-2-7B & 0.965 & 19 \\
    MinerU2.5-Pro (2604) & 0.925 & 40 \\
    MinerU2.5-Pro (2605) & 0.753 & 129 \\
    \bottomrule
  \end{tabular}

  \vspace{4pt}
  \begin{minipage}{0.95\linewidth}\footnotesize
  \textbf{Protocol note.} Measured on the sealed validation split. Findability asks a question of the transcript alone --- for each reference value, is it present at all, allowing the formatting variants a German invoice actually uses --- so it isolates reading from structuring. \emph{Corrected} means that 23 reference values which no channel could locate anywhere on the page, including the embedded text layer, are excluded as unanswerable rather than charged to every candidate equally. The embedded text layer is the ceiling: it is what the document itself contains, so no vision system can exceed it. This metric has a known blind spot --- a value can be present as a string while the surrounding structure is destroyed --- which is why it decided nothing on its own.
  \end{minipage}
\end{table}


The two finalists are statistically indistinguishable on corrected reading quality. An exact
McNemar test over the 963 paired per-cell findability outcomes --- 931 cells both readers
find, 16 only the general-purpose model finds, 13 only the transcription model finds, and
3 that neither finds, the same three cells the embedded text layer itself lacks ---
gives $p = 0.71$, so the observed ordering is well within what chance produces on discordant
counts this small.

One qualification on that test. It treats the 963 cells as independent trials, and they are
not: they nest within 29 invoices, and cells on the same page share a layout, a scan quality
and an issuer's typography, so outcomes within an invoice are correlated. The consequence runs
in the conservative direction here. Positive intra-invoice correlation reduces the effective
sample size, which \emph{widens} the interval a difference must clear and pushes $p$ further
from significance --- so a test that ignored the nesting cannot have manufactured the
null result it reports. A clustered alternative would be required to claim a difference from
these data; it is not required to decline to claim one, which is all this passage does.

What does separate the two readers is that they \emph{fail in different ways},
which is the substantive finding:

\begin{itemize}
  \item The document-transcription model produces near-flawless characters but silently drops
        letterhead and footer blocks --- which is exactly where a German invoice prints the
        tax identifier, the bank details and the seller's registered identity.
  \item The general-purpose model reads every margin block but makes occasional
        character-level slips \emph{inside} values: a transposed letter, a duplicated digit.
\end{itemize}

For this application the second failure mode is preferable, and the reason is the honest-null
contract rather than the score. A dropped block produces an absence, and under the honest-null
contract an absence is exactly what a genuinely missing field also produces --- so the failure
is camouflaged by the system's own convention, and it removes several legally mandatory
particulars at once, systematically, on every document sharing that layout. A character slip
produces a value that is present and wrong: it is attached to a named field, so a reviewer has
something concrete to check against the page, and the misreading is per-value rather than
per-region. An error a reviewer can see is worth more than an error of equal magnitude that
looks like a legitimate blank. The general-purpose model was selected.

\subsection{A pre-registered test that failed}

The obvious next question is whether the larger sibling of the selected model is better. Because
a favourable reading was available in advance --- more parameters, better reading --- the test
was pre-registered with a conjunctive rule: the challenger had to pass on reading quality,
\emph{and} introduce no new failure class, \emph{and} stay within a stated wall-clock factor.
All three, or the incumbent stands.

It passed on reading quality, with fewer real misses than the incumbent
(Table~\ref{tab:reader-findability}). It passed on wall-clock. It failed the middle criterion:
on one document it collapsed into a degenerate repetition loop, emitting thousands of repeated
characters and losing several fields entirely --- a failure mode the incumbent does not exhibit
on any document in the corpus.

A rare catastrophic failure is worse than a slightly higher rate of small ones for a system
whose value proposition is that a practitioner can trust its output. The incumbent stands. This
is reported at length because the conjunctive rule is what made the outcome defensible: had the
criteria been weighed after the fact, a marginally better reading score against a single
anomaly would have been easy to rationalise in either direction.

\section{Where the Shortfall Sits}
\label{sec:results-attribution}

The central question is which stage limits the pipeline. Three instruments answer it
(\S\ref{sec:method-attribution}), and their verdict is unambiguous.

Table~\ref{tab:attribution-clusters} and Figure~\ref{fig:attribution} give the same structuring
model on the reader's transcript and on text rendered from the reference. Only the input text
differs, so the gap between the two bars is attributable to reading. The gap is present in every
field cluster and it is large --- smallest on the core header fields, largest on the repeating
structures.

% GENERATED FILE -- DO NOT EDIT BY HAND.
% Regenerate with: make thesis-assets  (scripts/thesis_assets.py)
% Editing this file by hand breaks the guarantee that every number in the
% manuscript is derived from committed evidence (ADR-055).

% Source(s): data/finetune/attribution-val.json; data/finetune/attribution-oracle-val.json

\begin{table}[htbp]
  \centering
  \caption{Error attribution by field cluster. The gap column is the part of the shortfall that disappears when the reading stage is removed.}
  \label{tab:attribution-clusters}
  \begin{tabular}{lrrrrrr}
    \toprule
    Cluster & $n$ & Reader & Perfect text & Gap & FN\textsubscript{readable} & FN\textsubscript{unreadable} \\
    \midrule
    Core header fields & 29 & 0.8370 & 0.9750 & +0.1380 & 21 & 77 \\
    Extended header fields & 29 & 0.6873 & 0.9003 & +0.2130 & 43 & 39 \\
    Line items & 27 & 0.5549 & 0.9706 & +0.4157 & 82 & 157 \\
    Tax breakdown & 27 & 0.5903 & 0.9864 & +0.3961 & 47 & 31 \\
    Early-payment discount & 5 & 0.3333 & 0.7500 & +0.4167 & 8 & 0 \\
    \bottomrule
  \end{tabular}

  \vspace{4pt}
  \begin{minipage}{0.95\linewidth}\footnotesize
  \textbf{Protocol note.} Measured with a superseded reader (\texttt{granite-docling-258M}), which was replaced afterwards. These figures therefore establish \emph{where} the loss sits and \emph{which stage} causes it; they are not a system score and must not be read as one. Both score columns are cell-pooled within the cluster and come from the same structurer on the same documents --- only the text it was given differs, so the gap is attributable to reading. A miss is \emph{unreadable} when the expected value is absent from the transcript the structurer received: no structurer could have recovered it. It is \emph{readable} when the value was present and the structurer still failed to place it. That split is what makes the two stages separately accountable, and it is measured per miss rather than inferred from the totals.
  \end{minipage}
\end{table}


\begin{figure}[htbp]
  \centering
  \includegraphics[width=\textwidth]{attribution-by-cluster}
  \caption[Reader-limited versus structurer-limited performance]{The same structuring model on
  the reading model's transcript and on text rendered from the reference data, per field
  cluster. Only the input text differs between the two bars, so the annotated gap is the part
  of the shortfall attributable to reading. Measured with a superseded reading model: this
  attributes a shortfall between stages, it does not score the delivered system.}
  \label{fig:attribution}
\end{figure}

Table~\ref{tab:attribution-shares} attributes every individual signal-bearing error to a stage
by the per-miss test: was the value in the transcript at all? Roughly half of all errors are
values the transcript never contained, which no structuring model could have recovered. That
per-miss split is a \emph{lower bound} on the reading share, because it credits the structurer
with every value that was readable yet mangled in transit --- a transcript can contain a value
in a form so degraded that failing to place it is still, causally, a reading problem. The
pipeline-level decomposition is starker and uses explicit denominators: of the roughly 0.32
shortfall from a perfect score, the 0.28 between the two arms of
Table~\ref{tab:attribution-clusters} is attributable to reading, because only the input text
differs between them; the 0.04 the structurer leaves even on perfect text bounds its own
capability gap. By the per-miss count the stages split evenly; by the pipeline-level bound,
reading accounts for the substantial majority. Both statements are true, they answer different
questions, and the protocol note to Table~\ref{tab:attribution-shares} carries the exact
figures for each.

% GENERATED FILE -- DO NOT EDIT BY HAND.
% Regenerate with: make thesis-assets  (scripts/thesis_assets.py)
% Editing this file by hand breaks the guarantee that every number in the
% manuscript is derived from committed evidence (ADR-055).

% Source(s): data/finetune/attribution-val.json

\begin{table}[htbp]
  \centering
  \caption{Every signal-bearing error on the sealed split, attributed to the stage that caused it.}
  \label{tab:attribution-shares}
  \begin{tabular}{lrr}
    \toprule
    Cause & Errors & Share \\
    \midrule
    Reading: value absent from the transcript & 304 & 50.9\,\% \\
    Structuring: value present, not placed & 201 & 33.7\,\% \\
    Structuring: value emitted for an absent field & 92 & 15.4\,\% \\
    \midrule
    \textbf{Structuring, combined} & \textbf{293} & \textbf{49.1\,\%} \\
    \midrule
    \textbf{All signal-bearing errors} & \textbf{597} & \textbf{100.0\,\%} \\
    \bottomrule
  \end{tabular}

  \vspace{4pt}
  \begin{minipage}{0.95\linewidth}\footnotesize
  \textbf{Protocol note.} Same superseded-reader caveat as Table~\ref{tab:attribution-clusters}: this attributes a shortfall, it does not score a system. Signal-bearing means the cell contributes to F$_1$; correct absences are excluded, because counting them would make the split a function of how often fields are empty. On this corpus the pipeline scored 0.6771 against 0.9608 for the same structurer on perfect text, so of the 0.2837 gap the reading stage accounts for roughly 0.14 and the structurer for the remainder.
  \end{minipage}
\end{table}


\textbf{This measurement is diagnostic, not a system score, and the distinction is not
cosmetic.} It was taken with a reading model that was subsequently replaced --- the survey
reader, not the selected one. What it establishes is the \emph{attribution method} and the
\emph{ordering of the two stages}, both of which survive the reader change. What it does not
establish is any figure for the delivered system; those are in
\S\ref{sec:results-baseline} and \S\ref{sec:results-heldout}.

The finding is the thesis's central empirical claim: \emph{on this task, with local models at
this scale, reading is the binding constraint and structuring is not.} Chapter~\ref{ch:discussion}
takes up what follows from that for system design.

\section{The Sealed Baseline}
\label{sec:results-baseline}

Table~\ref{tab:sealed-val-arms} gives the zero-shot pipeline on the sealed validation split, in
both numeric precisions and on both input conditions. Four observations.

% GENERATED FILE -- DO NOT EDIT BY HAND.
% Regenerate with: make thesis-assets  (scripts/thesis_assets.py)
% Editing this file by hand breaks the guarantee that every number in the
% manuscript is derived from committed evidence (ADR-055).

% Source(s): data/finetune/eval-zeroshot-bf16-val.json; data/finetune/eval-oracle-bf16-val.json; data/finetune/eval-zeroshot-qwen-adr059-val.json; data/finetune/eval-oracle-adr059-val.json

\begin{table}[htbp]
  \centering
  \caption{Zero-shot structurer performance on the sealed validation split, in both numeric precisions and on both input conditions.}
  \label{tab:sealed-val-arms}
  \begin{tabular}{lrrr}
    \toprule
    Arm & Overall F$_1$ & Flat F$_1$ & Spurious rate \\
    \midrule
    bf16 / CUDA \quad reader transcript & 0.8480 & 0.8843 & 0.1575 \\
    bf16 / CUDA \quad perfect text & 0.9778 & 0.9832 & 0.0204 \\
    \midrule
    4-bit / Apple silicon \quad reader transcript & 0.8257 & 0.8649 & 0.1462 \\
    4-bit / Apple silicon \quad perfect text & 0.9719 & 0.9743 & 0.0145 \\
    \bottomrule
  \end{tabular}

  \vspace{4pt}
  \begin{minipage}{0.95\linewidth}\footnotesize
  \textbf{Protocol note.} Corpus: the 29 sealed synthetic invoices, never used for fitting or selection. Structurer: \texttt{google/gemma-4-E4B-it}, greedy decoding, 29/29 parsed in every arm. Ground truth is extracted from each invoice's own embedded XML and is exact by construction. \emph{Reader transcript} is the deployed pipeline: the structurer reads what the vision model transcribed. \emph{Perfect text} replaces that transcript with text rendered from the ground truth, so it measures the structurer with the reading stage removed --- an instrument, not a deliverable. Both aggregations are means over the 29 per-invoice scores. Overall pools flat fields with repeating-group cells; flat covers the header fields only. Spurious rate is the share of emitted values for fields the document does not carry.
  \end{minipage}
\end{table}


\textbf{The deployable configuration --- 4-bit weights on a laptop --- reaches a score in the
mid-eighties on flat header fields.} That is the figure a practitioner would experience, on
generated pages.

\textbf{The structuring stage is near its ceiling on perfect text.} Given text rendered from the
reference, the same model scores in the high nineties. The structuring stage is not what limits
this pipeline, which corroborates \S\ref{sec:results-attribution} on the current reader rather
than the superseded one.

\textbf{Full numeric precision is worth a measurable amount} --- the difference between the two
stacks is small but not negligible, and \S\ref{sec:results-adaptation} shows why that
seemingly minor quantity mattered enormously.

\textbf{The spurious-emission rate is markedly lower on perfect text than on transcripts.} The
model invents fewer values when the input is clean. This is intuitive and it is also the
mechanism that the adaptation result turns on.

\section{Adaptation: A Negative Result}
\label{sec:results-adaptation}

The scope decision for this thesis made a fine-tuning experiment conditional: it would be
attempted only if the zero-shot pipeline fell short of a stated threshold, and only after
cheaper explanations had been eliminated. Both conditions were met, so the experiment ran.

\subsection{Cheaper explanations were eliminated first}

Before any training, every underperforming field was classified by cause: value absent from the
input, value present but never named in the field specification, or value present and named but
rejected by a comparison rule. A gap of the second kind is a specification defect, and
repairing it by fine-tuning would credit the adapter with a gain that a one-line change would
have produced.

The classification escalated three candidate fields for per-document inspection.
\emph{All three turned out not to be specification-fixable.} One was a net-versus-gross confusion between two
labels that were both already grounded, and it landed mostly on the \emph{perfect-text} arm ---
which means there was no headroom to recover. One emitted a character-identical value on almost
every document that carried it, with most of its remaining misses caused by a
normalisation asymmetry, not by the model. One was corrupted reference data on two of its
three misses. The verdict was zero warranted specification repairs, so the fine-tuning
experiment was not being handed a gap that belonged to the prompt.

\subsection{The design: two adapters, one variable}

Rather than train one adapter and measure it, two were trained, differing only in the input
distribution they saw: one on the reader's actual transcripts, one on text rendered from the
reference. Both are low-rank adapters \parencite{hu2021lora} on the structuring model's language
tower only --- its vision and audio components stay frozen and receive no input, consistent
with how the pipeline uses the model (\S\ref{sec:method-architecture}). Each was then
evaluated on both conditions, giving the grid in Table~\ref{tab:finetune-grid} and
Figure~\ref{fig:finetune-grid}.

The reason for two is that a single measurement would have been uninterpretable. An adapter
trained and evaluated on transcripts conflates two distinct claims --- learning the output
schema, and learning to survive reading noise --- and a gain on one is indistinguishable from a
gain on the other. The grid separates them. The perfect-text adapter is an instrument and is
never proposed as a deliverable.

The recipe, in Table~\ref{tab:hyperparameters}, was registered before either run, along with the
selection rule, the requirement of a matched-precision baseline, and the commitment to publish a
null result. No hyperparameter sweep was performed, which is a limitation and is recorded as
one in Chapter~\ref{ch:limitations-future-work}.

% GENERATED FILE -- DO NOT EDIT BY HAND.
% Regenerate with: make thesis-assets  (scripts/thesis_assets.py)
% Editing this file by hand breaks the guarantee that every number in the
% manuscript is derived from committed evidence (ADR-055).

% Source(s): data/finetune/adapter/horus_training_provenance.json; data/finetune/adapter-oracle/horus_training_provenance.json

\begin{table}[htbp]
  \centering
  \caption{The adaptation recipe, as recorded by the training runs themselves.}
  \label{tab:hyperparameters}
  \begin{tabular}{lrr}
    \toprule
    Setting & Reader arm & Perfect-text arm \\
    \midrule
    Rank & 8 & identical \\
    LoRA alpha & 16.0 & identical \\
    Dropout & 0.05 & identical \\
    Learning rate & 0.0001 & identical \\
    Schedule & cosine & identical \\
    Warm-up fraction & 0.03 & identical \\
    Epoch budget & 6.0 & identical \\
    Batch size & 1 & identical \\
    Gradient accumulation & 8 & identical \\
    Sequence budget (auto-sized) & 6144 & 3840 \\
    Seed & 42 & identical \\
    Fitting examples & 100 & 100 \\
    Adapted modules & 258 & 258 \\
    Dev-slice size & 17 & 17 \\
    \bottomrule
  \end{tabular}

  \vspace{4pt}
  \begin{minipage}{0.95\linewidth}\footnotesize
  \textbf{Protocol note.} Read from the provenance file each run wrote, not from the configuration it was asked to use, so the table reports what executed. The sequence budget is the one setting that differs, and it is derived rather than chosen: the trainer sizes it to the longest example in that arm's own training set, rounded up to a 256-token boundary and capped at 8192. Neither cap binds, so no training example was truncated in either arm. The difference is therefore a consequence of the variable under study --- reader transcripts are longer than rendered text --- and not an independent second variable. Adapted modules are confined to the text tower; the vision tower stays frozen throughout.
  \end{minipage}
\end{table}


% GENERATED FILE -- DO NOT EDIT BY HAND.
% Regenerate with: make thesis-assets  (scripts/thesis_assets.py)
% Editing this file by hand breaks the guarantee that every number in the
% manuscript is derived from committed evidence (ADR-055).

% Source(s): data/finetune/eval-zeroshot-bf16-val.json; data/finetune/eval-oracle-bf16-val.json; data/finetune/eval-ft-reader-on-reader-val.json; data/finetune/eval-ft-reader-on-oracle-val.json; data/finetune/eval-ft-oracle-on-reader-val.json; data/finetune/eval-ft-oracle-on-oracle-val.json

\begin{table}[htbp]
  \centering
  \caption{The 2$\times$2 adaptation grid: two adapters, each evaluated on both input conditions. All four cells regress against the matched-stack baseline.}
  \label{tab:finetune-grid}
  \begin{tabular}{llrrr}
    \toprule
    Model & Evaluated on & Overall F$_1$ & $\Delta$ vs baseline & Spurious rate \\
    \midrule
    \emph{no adapter} (baseline) & reader transcript & 0.8480 & --- & 0.1575 \\
    \emph{no adapter} (baseline) & perfect text & 0.9778 & --- & 0.0204 \\
    \midrule
    adapter trained on reader transcripts & reader transcript & 0.8246 & -0.0234 & 0.2012 \\
    adapter trained on reader transcripts & perfect text & 0.9583 & -0.0196 & 0.0456 \\
    \midrule
    adapter trained on perfect text & reader transcript & 0.8354 & -0.0126 & 0.1692 \\
    adapter trained on perfect text & perfect text & 0.9303 & -0.0476 & 0.0287 \\
    \bottomrule
  \end{tabular}

  \vspace{4pt}
  \begin{minipage}{0.95\linewidth}\footnotesize
  \textbf{Protocol note.} Corpus and structurer as in Table~\ref{tab:sealed-val-arms}. Every arm here is bf16 on CUDA, including the baseline: comparing a bf16 adapter against the 4-bit deployment baseline would measure the adapter \emph{and} a precision change at once. The two adapters differ only in the input distribution they were trained on; rank, scaling, dropout, schedule, seed and target modules are identical, and both selected the same checkpoint. Deltas are against the baseline in the same evaluation condition, never across conditions. The adapter trained on perfect text is an instrument for attributing the shortfall, not a candidate for deployment.
  \end{minipage}
\end{table}


\begin{figure}[htbp]
  \centering
  \includegraphics[width=\textwidth]{finetune-grid}
  \caption[The adaptation grid]{Two adapters, each evaluated on both input conditions, against
  the matched-precision baseline. Every cell regresses. Note that the baseline shown is
  re-measured in the adapters' own numeric precision; the deployed 4-bit baseline is lower, and
  comparing against it instead is what Table~\ref{tab:precision-confound} shows to be
  misleading.}
  \label{fig:finetune-grid}
\end{figure}

\subsection{The result}

\textbf{All four cells of the grid regress.} Table~\ref{tab:finetune-grid} gives the deltas
against the matched baseline in each evaluation condition. No arrangement of the two adapters
and the two conditions produces an improvement over doing nothing.

Two further observations make the result interpretable instead of merely negative.

\textbf{The spurious-emission rate rises where the score falls.} The adapter trained on
transcripts has the highest spurious-emission rate of any arm measured in this thesis
(Table~\ref{tab:finetune-grid}) --- markedly above its own baseline. The adapter did not lose
the ability to place values; it gained a propensity to emit them where the document carries
none. That is a specific, mechanistic finding, not a general degradation, and
Chapter~\ref{ch:discussion} reads it against the published literature on the same effect.

\textbf{Both arms overfit after one epoch.} Table~\ref{tab:devloss} and
Figure~\ref{fig:devloss} give the loss on the held-back slice by epoch. Both arms reach their
minimum after a single epoch and never return to it within the six-epoch budget. The selection
rule --- lowest loss on that slice --- therefore chose the earliest checkpoint in both arms,
which is not a tuning decision but the rule doing what it was registered to do. The six-epoch
budget existed to make the turn visible, not to be used.

% GENERATED FILE -- DO NOT EDIT BY HAND.
% Regenerate with: make thesis-assets  (scripts/thesis_assets.py)
% Editing this file by hand breaks the guarantee that every number in the
% manuscript is derived from committed evidence (ADR-055).

% Source(s): data/finetune/adapter/horus_training_provenance.json; data/finetune/adapter-oracle/horus_training_provenance.json

\begin{table}[htbp]
  \centering
  \caption{Held-back dev-slice loss by epoch for both adapters. Both minimise after one epoch and never recover.}
  \label{tab:devloss}
  \begin{tabular}{rrrl}
    \toprule
    Epoch & Reader arm & Perfect-text arm &  \\
    \midrule
    1 & 0.1956 & 0.0965 & \quad$\leftarrow$ selected \\
    2 & 0.3918 & 0.3095 &  \\
    3 & 0.4438 & 0.3799 &  \\
    4 & 0.4183 & 0.3644 &  \\
    5 & 0.3857 & 0.3150 &  \\
    6 & 0.3801 & 0.3078 &  \\
    \bottomrule
  \end{tabular}

  \vspace{4pt}
  \begin{minipage}{0.95\linewidth}\footnotesize
  \textbf{Protocol note.} Loss is measured on a 17-invoice dev slice carved out of the training split, disjoint from both the fitting set and the sealed validation split, so selecting on it leaks nothing into the reported score. The two arms share the same dev slice. The selection rule --- lowest dev loss --- was fixed before the runs and is enforced by the trainer rather than applied by hand. Loss rises by 2.00$\times$ (reader) and 3.21$\times$ (perfect text) at the second epoch and stays above the epoch-one value for the remaining budget, so the six-epoch budget was never a target: it existed to make the turn visible.
  \end{minipage}
\end{table}


\begin{figure}[htbp]
  \centering
  \includegraphics[width=0.92\textwidth]{devloss-curve}
  \caption[Dev-slice loss by epoch]{Loss on the held-back development slice, by epoch, for both
  adapters. Both minimise after one epoch; the registered selection rule therefore chose the
  earliest checkpoint in both arms. The slice is carved from the fitting set and is disjoint
  from the sealed validation split, so selecting on it leaks nothing into any reported score.}
  \label{fig:devloss}
\end{figure}

With 100 fitting examples against a model of this size, immediate overfitting is not surprising.
What the grid adds is that the overfitting is not the whole story: the perfect-text adapter,
evaluated on perfect text, also regresses, and it saw clean input throughout.

\subsection{The comparison that would have inverted the conclusion}

Table~\ref{tab:precision-confound} is the most important table in this chapter, and it is not
about the adapter.

The adapters were trained in full precision on target-class hardware because they exceeded the
floor's memory. The deployed baseline runs in 4-bit precision at the floor. Comparing the two
directly measures the adapter and a precision change simultaneously --- and because the
adapter's damage and full precision's benefit are of similar size, they nearly cancel. The
unmatched comparison reads as no measurable harm. The matched comparison, against a baseline
re-measured in the adapter's own precision, shows a regression roughly twenty times larger.

% GENERATED FILE -- DO NOT EDIT BY HAND.
% Regenerate with: make thesis-assets  (scripts/thesis_assets.py)
% Editing this file by hand breaks the guarantee that every number in the
% manuscript is derived from committed evidence (ADR-055).

% Source(s): data/finetune/eval-ft-reader-on-reader-val.json; data/finetune/eval-zeroshot-bf16-val.json; data/finetune/eval-zeroshot-qwen-adr059-val.json

\begin{table}[htbp]
  \centering
  \caption{Why the baseline had to be re-measured in the adapter's own numeric precision. The unmatched comparison understates the regression by a factor of 21.}
  \label{tab:precision-confound}
  \begin{tabular}{lrrrl}
    \toprule
    Comparison & Adapter & Baseline & $\Delta$ & Reads as \\
    \midrule
    unmatched (adapter bf16 vs baseline 4-bit) & 0.8246 & 0.8257 & -0.0011 & ``no measurable harm'' \\
    matched (both bf16) & 0.8246 & 0.8480 & -0.0234 & \textbf{a real regression} \\
    \bottomrule
  \end{tabular}

  \vspace{4pt}
  \begin{minipage}{0.95\linewidth}\footnotesize
  \textbf{Protocol note.} Both rows use the same adapter and the same sealed corpus; only the baseline changes. The deployed baseline runs 4-bit on Apple silicon, the adapter runs bf16 on CUDA, and full precision is worth +0.0223 on this corpus. Because the adapter costs roughly what the precision gain is worth, the two nearly cancel: the unmatched difference is small enough to read as neutral. The regression is only visible once the baseline is re-measured in the adapter's own precision, which is why that re-measurement was made a precondition of the study rather than an afterthought.
  \end{minipage}
\end{table}


Had the baseline not been re-measured, this thesis would have reported that a fine-tuned
structurer performs about the same as an untuned one, and would have drawn the wrong conclusion
from a correctly computed number. Chapter~\ref{ch:measurement-validity} treats the general
lesson; here the point is narrower and sharper: \emph{the conclusion of this experiment is a
property of the experimental design, not of the measurement precision.}

\section{Real Invoices}
\label{sec:results-heldout}

The synthetic corpus answers whether the pipeline can read a well-formed generated page. The
held-out corpus answers whether that survives real documents. This is the result the thesis
exists to produce.

\subsection{The headline}

Table~\ref{tab:heldout-headline} gives header-field extraction on 39 real invoices, zero-shot,
with open-weights models only, against an answer key adjudicated across three independent
channels and signed off document by document. All 39 documents parsed and scored; none was
dropped.

\textbf{Venue.} Every model involved is open-weights and nothing on the inference path makes a
network call. The reading pass ran at full precision and native page resolution on
target-class hardware --- a single 24\,\ac{GB} graphics processor, which transcribed all
58 pages in roughly sixteen minutes --- and the structuring pass ran at the floor, on the
laptop in its deployed 4-bit precision (\S\ref{sec:method-repro}). The figure below is
therefore the pipeline's accuracy on the hardware class a company would buy for it. What the
floor configuration's capped reader would score on the same documents is a separate, smaller
question, recorded as open in \S\ref{sec:lim-hardware}.

% GENERATED FILE -- DO NOT EDIT BY HAND.
% Regenerate with: make thesis-assets  (scripts/thesis_assets.py)
% Editing this file by hand breaks the guarantee that every number in the
% manuscript is derived from committed evidence (ADR-055).

% Source(s): eval/heldout-breakdown.json; data/self-collected/_eval/eval-zeroshot-heldout-adr065.json (private)

\begin{table}[htbp]
  \centering
  \caption{Header-field extraction on the held-out corpus of real invoices, zero-shot, with open-weights models only.}
  \label{tab:heldout-headline}
  \begin{tabular}{lr}
    \toprule
    Quantity & Value \\
    \midrule
    Documents scored & 39 / 39 \\
    Mean per-invoice F$_1$ & 0.8825 \\
    Cell-pooled F$_1$ & 0.8987 \\
    Precision & 0.9530 \\
    Recall & 0.8503 \\
    \midrule
    True positives & 568 \\
    False positives & 28 \\
    False negatives & 100 \\
    \bottomrule
  \end{tabular}

  \vspace{4pt}
  \begin{minipage}{0.95\linewidth}\footnotesize
  \textbf{Protocol note.} Corpus: real invoices collected from the author's own correspondence, never used for fitting, prompt design or selection. Answer key: the signed-off key, adjudicated across three independent reading channels and confirmed document by document by the author. Scope: the 34 registered header fields; repeating groups are excluded structurally, because their rows were never author-reviewed. Precision above recall means roughly four errors in five are a field left empty rather than a field invented --- for an accounting tool, the safer direction to fail in, because a gap is visible to a reviewer and a fabrication is not. \textbf{Venue.} Every model here is open-weights and no inference path makes a network call. Reading ran at full precision and native resolution on target-class hardware (a single 24\,GB graphics processor); structuring ran at the floor, on a 16\,GB laptop in its deployed 4-bit precision (\S\ref{sec:method-repro}). The figure is the pipeline's accuracy on the hardware class a firm would buy for it; the floor's capped-resolution reading is not measured on this corpus (\S\ref{sec:lim-hardware}). An earlier ruler scored the same stored generations at 0.8767; the difference is the instrument, not the system, and is accounted for in Chapter~\ref{ch:measurement-validity}.
  \end{minipage}
\end{table}


The result is that performance largely survives the transition from generated pages to real
ones. That is not a trivial finding: the documents were produced by other people's billing
systems, in two languages, with layouts the pipeline was never shown, and roughly a quarter of
them were phone-scanned paper rather than exported files.

\textbf{The headline pools two languages, and the research question names one.} Of the 39
documents, 28 are German and 11 are English, so the figure above is not a German-only
measurement. Summing the two German rows of Table~\ref{tab:heldout-by-channel} gives 417 true
positives against 25 false positives and 84 false negatives, i.e.\ a pooled F$_1$ of
$0.884$ and a per-invoice mean of $0.862$ over the 28 German invoices --- both slightly
\emph{below} the mixed-corpus headline, because the English email subset is the
best-read channel in the corpus (Table~\ref{tab:heldout-by-channel}). Pooling therefore
flatters the German figure rather than the reverse, and the German-specific claim is the
weaker of the two. The English documents are retained in the headline because they were
collected under the same procedure and adjudicated under the same rule, and excluding them
after seeing the per-channel split would be the selection this thesis spends
Chapter~\ref{ch:measurement-validity} arguing against.

\subsection{The error profile matters more than the score}

Precision is substantially higher than recall (Table~\ref{tab:heldout-headline}). Roughly four
errors in five are a field left empty, not a field invented.

For an accounting application this is the safer direction to fail in, and the asymmetry is not
incidental --- it is the honest-null contract of \S\ref{sec:design-honest-null}, with the
gap-versus-fabrication argument made there, showing up in the numbers. A system with the
opposite profile at the same F$_1$ would be considerably less useful, which is why the
spurious-emission rate is tracked as a first-class metric throughout instead of folded into
an aggregate.

\subsection{Degraded capture is the real cost}

Table~\ref{tab:heldout-by-channel} and Figure~\ref{fig:heldout-channel} break the corpus down by
language and acquisition channel, and this is the most useful measurement in the thesis for
anyone who might deploy the system.

\textbf{Language costs little. Capture channel costs a great deal.} English and German
email-native documents score close together. Phone-scanned German documents fall away sharply
--- a gap of more than eleven points of mean per-invoice F$_1$ against the email-native
documents. And the scans in this corpus are the \emph{favourable} kind of degraded input:
scanning-app output that has already been deskewed and contrast-enhanced
(\S\ref{sec:method-data}), not raw camera photographs. The gap for uncorrected captures
should be expected to be larger, not smaller.

Two properties of that gap deserve emphasis.

\textbf{It survives the ruler correction.} The gap is present under both the superseded and the
current scoring rules, at similar magnitude. It is not an artefact of the answer-key work
described in Chapter~\ref{ch:measurement-validity}.

\textbf{It is a recall gap, not a precision gap.} Figure~\ref{fig:heldout-pr} separates the two:
precision holds up on the scanned documents while recall falls. \textbf{Degraded input makes the
system abstain, not hallucinate.} That is the behaviour one would want if forced to choose, and it is
worth noting that nothing in the design explicitly asked for it.

% GENERATED FILE -- DO NOT EDIT BY HAND.
% Regenerate with: make thesis-assets  (scripts/thesis_assets.py)
% Editing this file by hand breaks the guarantee that every number in the
% manuscript is derived from committed evidence (ADR-055).

% Source(s): eval/heldout-breakdown.json; data/self-collected/_eval/eval-zeroshot-heldout-adr065.json (private)

\begin{table}[htbp]
  \centering
  \caption{Held-out performance by language and capture channel. Photographed documents cost far more than language does.}
  \label{tab:heldout-by-channel}
  \begin{tabular}{lrrrrr}
    \toprule
    Channel & $n$ & Mean/invoice & Pooled & Precision & Recall \\
    \midrule
    English, email-native PDF & 11 & 0.9345 & 0.9408 & 0.9805 & 0.9042 \\
    German, email-native PDF & 18 & 0.9027 & 0.9079 & 0.9533 & 0.8667 \\
    German, phone scan & 10 & 0.7889 & 0.8371 & 0.9225 & 0.7661 \\
    \midrule
    \textbf{All documents} & \textbf{39} & \textbf{0.8825} & \textbf{0.8987} & \textbf{0.9530} & \textbf{0.8503} \\
    \bottomrule
  \end{tabular}

  \vspace{4pt}
  \begin{minipage}{0.95\linewidth}\footnotesize
  \textbf{Protocol note.} Both aggregations are derived from the same scoring run under the same ruler, so the row totals and the whole-corpus row are mutually consistent. Mean per invoice counts every document once regardless of how many fields it carries, and answers how the system behaves on a document handed to it. Cell-pooled sums all outcomes before computing one score, so field-dense documents pull harder, and answers what share of extracted cells is correct. Neither is a maximum over documents. Precision holds up on photographs while recall falls, meaning degraded input makes the system abstain rather than invent. Email-native documents average 0.9148 against 0.7889 for photographs, a gap of 12.6 points.
  \end{minipage}
\end{table}


\begin{figure}[htbp]
  \centering
  \includegraphics[width=\textwidth]{heldout-by-channel}
  \caption[Held-out performance by capture channel]{Header-field extraction on real invoices,
  by language and acquisition channel, under both aggregations. The dashed line is the
  whole-corpus mean. Language costs little; phone-scanned capture costs a great deal.}
  \label{fig:heldout-channel}
\end{figure}

\begin{figure}[htbp]
  \centering
  \includegraphics[width=\textwidth]{heldout-precision-recall}
  \caption[Precision and recall by capture channel]{The same corpus split into precision and
  recall. Precision holds up on phone-scanned documents while recall falls, which means degraded
  input makes the system abstain rather than invent.}
  \label{fig:heldout-pr}
\end{figure}

\textbf{This result is independently corroborated.} A benchmark published during this project,
covering 3{,}400 documents in 17 languages, reports an average drop of 17.8\,\% on photographed
versus digital documents for open-source parsing models \parencite{mdpbench2026}. Its degraded
condition is harsher than this corpus's --- printed pages photographed under controlled
adverse conditions, against scan-app output here --- so the two numbers are not directly
comparable; what they share is the direction and the order of magnitude. That agreement
matters because a single-annotator corpus of 39 documents cannot carry a general claim on its
own, and here it does not have to.

\subsection{What this result does not cover}

The held-out claim is about header fields. Repeating-group rows are excluded structurally,
because their reference rows were never author-reviewed and could not be reviewed affordably
(\S\ref{sec:method-heldout-gt}). Line-item extraction is measured on the synthetic corpus,
where the reference is exact and free, and is reported in
Table~\ref{tab:attribution-clusters}. The limitation, and the known route to closing it, are in
Chapter~\ref{ch:limitations-future-work}.

\section{Summary of Findings}
\label{sec:results-summary}

\begin{enumerate}
  \item \textbf{Reading is the binding constraint.} The same structuring model given text
        rendered from the reference scores far higher than it does on the reading model's
        transcript, in every field cluster. Roughly half of all individual errors are values
        the transcript never contained.
  \item \textbf{Local, zero-shot header extraction on real invoices works better than expected
        and fails in the right direction.} Precision well above recall; four errors in five are
        omissions, not fabrications.
  \item \textbf{Degraded capture, not language, is the dominant real-world cost}, at more than
        eleven points of mean per-invoice F$_1$, corroborated independently by a contemporary
        multilingual benchmark.
  \item \textbf{Parameter-efficient adaptation of the structuring stage does not help}, in any
        of four configurations, and it specifically increases the rate at which the model emits
        values for fields the document does not carry.
  \item \textbf{Two selection decisions could not be made on aggregate score.} Reader selection
        required per-miss adjudication because the reading-quality ordering inverted when the
        ruler was corrected; the adaptation conclusion required a precision-matched baseline,
        without which it would have been the opposite.
\end{enumerate}
